
\documentclass[aps,preprint]{revtex4}
%%%%%%%%%%%%%%%%%%%%%%%%%%%%%%%%%%%%%%%%%%%%%%%%%%%%%%%%%%%%%%%%%%%%%%%%%%%%%%%%%%%%%%%%%%%%%%%%%%%%%%%%%%%%%%%%%%%%%%%%%%%%
\usepackage{amsmath}
\usepackage{amssymb}

\setcounter{MaxMatrixCols}{10}
%TCIDATA{OutputFilter=LATEX.DLL}
%TCIDATA{Version=4.00.0.2321}
%TCIDATA{Created=Friday, April 05, 2002 10:30:19}
%TCIDATA{LastRevised=Tuesday, July 09, 2002 12:07:39}
%TCIDATA{<META NAME="GraphicsSave" CONTENT="32">}
%TCIDATA{<META NAME="DocumentShell" CONTENT="Articles\SW\REVTeX 4">}
%TCIDATA{Language=American English}
%TCIDATA{CSTFile=revtex4.cst}

\newtheorem{theorem}{Theorem}
\newtheorem{acknowledgement}[theorem]{Acknowledgement}
\newtheorem{algorithm}[theorem]{Algorithm}
\newtheorem{axiom}[theorem]{Axiom}
\newtheorem{claim}[theorem]{Claim}
\newtheorem{conclusion}[theorem]{Conclusion}
\newtheorem{condition}[theorem]{Condition}
\newtheorem{conjecture}[theorem]{Conjecture}
\newtheorem{corollary}[theorem]{Corollary}
\newtheorem{criterion}[theorem]{Criterion}
\newtheorem{definition}[theorem]{Definition}
\newtheorem{example}[theorem]{Example}
\newtheorem{exercise}[theorem]{Exercise}
\newtheorem{lemma}[theorem]{Lemma}
\newtheorem{notation}[theorem]{Notation}
\newtheorem{problem}[theorem]{Problem}
\newtheorem{proposition}[theorem]{Proposition}
\newtheorem{remark}[theorem]{Remark}
\newtheorem{solution}[theorem]{Solution}
\newtheorem{summary}[theorem]{Summary}
\newenvironment{proof}[1][Proof]{\noindent\textbf{#1.} }{\ \rule{0.5em}{0.5em}}
\input{tcilatex}

\begin{document}

\preprint{}
\title[Planned]{Long title that appears on the title page}
\author{Brian Weiner}
\affiliation{Physics, PSU}
\author{}
\author{}
\keywords{one two three}
\pacs{PACS number}

\begin{abstract}
Shell document for REV\TeX{} 4.
\end{abstract}

\volumeyear{year}
\volumenumber{number}
\issuenumber{number}
\eid{identifier}
\date[Date text]{date}
\received[Received text]{date}
\revised[Revised text]{date}
\accepted[Accepted text]{date}
\published[Published text]{date}
\startpage{1}
\endpage{}
\maketitle
\tableofcontents

\section{Outline of Coherent State}

\subsection{Resolutions of the Identity}

\subsubsection{Coherent States}

Resolutions of the identity can be based on families of Coherent States (CS)
as 
\begin{equation}
I_{N}=\int \left\vert \eta \right\rangle \left\langle \eta \right\vert d\mu
\left( \eta \right) ;\quad \left\vert \eta \right\rangle \in \mathcal{H}^{N}
\label{res_id}
\end{equation}%
where $\left\{ \left\vert \eta \right\rangle \right\} $ is an overcomplete
basis of $\mathcal{H}^{N}$ 
\begin{equation}
\eta :\mathrm{X\rightarrow }\mathbb{C}
\end{equation}%
$\mathrm{X}$ is the index set of the label space $\left\{ \eta \right\} $,
e.g. 
\begin{equation}
\mathrm{X}=\mathbb{N}
\end{equation}%
and 
\begin{equation}
\eta \equiv \left\{ \eta _{i}|1\leq i\leq r\right\} \equiv \left( \eta
_{1}\cdots \eta _{r}\right) \equiv \mathbf{\eta \in }\mathbb{C}^{r}
\end{equation}%
or $\mathrm{X}$ could be a more general measure space and $\left\{ \eta
\right\} $ some space of functions defined on $\mathrm{X}$. With each family
of CS's one can associate the Reproducing Kernel function (in Quantum
Chemistry called the overlap matrix) as 
\begin{equation}
\mathcal{K}\left( \eta ,\eta ^{\prime }\right) =\left\langle \eta |\eta
^{\prime }\right\rangle .
\end{equation}%
One special case of this 
\begin{equation}
\mathrm{X}=\mathbb{R}^{3N}
\end{equation}%
and 
\begin{equation}
\eta \left( \mathbf{x}_{1}\cdots \mathbf{x}_{N}\right) \equiv \eta \left( 
\mathsf{X}^{N}\right) =\left\langle \mathsf{X}^{N}|\eta \right\rangle ,\quad
\left\vert \eta \right\rangle \in \mathcal{H}^{N},\quad \eta \left( \mathbf{x%
}_{1}\cdots \mathbf{x}_{N}\right) \in L_{2}\left( \mathbb{R}^{3N}\right)
\end{equation}%
Another special case is possibly 
\begin{equation}
\mathrm{X}=\mathbb{R}^{3}
\end{equation}%
with 
\begin{equation}
\left\{ \eta \right\} \equiv L_{1+}\left( \mathbb{R}^{3}\right) ;\quad \int
\eta \left( \mathbf{x}_{1}\right) d\mathbf{x}_{1}=N
\end{equation}%
By Harrimans theorem every density corresponds to an IPS\ state, but not
uniquely i.e. a given density can be associated with a class of IPS's, so
the question to be decided is whether an integral over class representatives
is sufficient to produce a resolution of the identity.

If the CS's are group CS's then 
\begin{equation}
I_{N}=\int \mathcal{R}\left( g\left( \eta \right) \right) \left\vert \eta
_{ref}\right\rangle \left\langle \eta _{ref}\right\vert \mathcal{R}\left(
g\left( \eta \right) \right) ^{\dagger }d\mu \left( \eta \right)
\label{groupres}
\end{equation}%
where $\mathcal{R}$ is a representation of a group $\mathfrak{G}$ acting in $%
\left\{ \mathcal{H}^{N}\right\} $, where the group elements are
parameterized by $\eta $. If $\mathcal{R}$ is irreducible then eq. $\left( %
\ref{res_id}\right) $ follows. While if $\mathcal{R}$ is reducible the
integral does not in general produce the identity operator, but rather the
operator 
\begin{equation}
J=\sum_{1\leq \kappa \leq m}\alpha \left( \kappa \right) P_{\kappa }
\end{equation}%
where $m$ is the number of irreps (not counting degeneracy) and $P_{\kappa }$
is a projector onto invariant subspaces that carry irreps. In the $\infty $
dimensional case there may not be any finite dimensional invariant subspaces
even if $\mathcal{R}$ is reducible, and even if $\mathfrak{G}$ is abelian,
so eq. $\left( \ref{groupres}\right) $ could then be 
\begin{equation}
J=\sum_{1\leq \kappa \leq m}\alpha \left( \kappa \right) P_{\kappa }
\end{equation}%
where $\left\{ P_{\kappa }\right\} $ are all projectors onto invariant $%
\infty $ dimensional subspaces.

\subsubsection{Generalized States}

Another set of resolutions of the identity can be based on \emph{generalized 
}states as 
\begin{equation}
I_{N}=\int \left\vert \mathsf{z}^{N}\right\rangle \left\langle \mathsf{z}%
^{N}\right\vert d\mathsf{z}^{N}
\end{equation}%
where $\left\vert \mathsf{z}^{N}\right\rangle \notin \mathcal{H}^{N}$, $%
\mathsf{z}^{N}\in \mathrm{X}$ a linear atomic measure space, and form an
overcomplete orthonormal basis of $\mathcal{H}^{N}$ i.e. 
\begin{equation}
\left\langle \mathsf{z}^{N}|\mathsf{z}^{\prime N}\right\rangle =\delta
\left( \mathsf{z}^{N}-\mathsf{z}^{\prime N}\right)
\end{equation}%
which also defines the amplitudes, $\delta _{\mathsf{z}^{\prime N}}\left( 
\mathsf{z}^{N}\right) \equiv \delta \left( \mathsf{z}^{N}-\mathsf{z}^{\prime
N}\right) $, of these basis functions. One could express this resolution of
the identity in "group" CS\ fashion as 
\begin{equation}
I_{N}=\int \breve{T}\left( \mathsf{z}^{N}\right) \left\vert \mathsf{z}%
_{ref}^{N}\right\rangle \left\langle \mathsf{z}_{ref}^{N}\right\vert \breve{T%
}\left( \mathsf{z}^{N}\right) ^{\dagger }d\mathsf{z}^{N}
\end{equation}%
where $\left\vert \mathsf{x}^{N}\right\rangle =\breve{T}\left( \mathsf{x}%
^{N}\right) \left\vert \mathsf{z}_{ref}^{N}\right\rangle $ i.e. 
\begin{eqnarray}
\delta _{\mathsf{x}^{N}} &=&\delta _{\breve{T}\left( \mathsf{x}^{N}\right) 
\mathsf{z}_{ref}^{N}}  \notag \\
\delta \left( \mathsf{x}^{N}-\mathsf{z}^{N}\right) &=&\delta \left( T\left( 
\mathsf{x}^{N}\right) \mathsf{z}_{ref}^{N}-\mathsf{z}^{N}\right)
\end{eqnarray}%
so that $\breve{T}\left( \mathsf{z}^{N}\right) $ is a representation of the
translational group of $\mathsf{Z}^{N}$.

The big difference between the generalized and CS resolution of the
identities is that the constraint of orthonormality for an overcomplete
basis cannot be fulfilled if the basis elements belong to the Hilbert space.

\subsection{ Coherent States based on Generalized States}

\subsection{Generation of States by Diagonal Operators}

If we consider the coordinate representation, all states (unnormalized) in $%
\mathcal{H}^{N}$ can be produced in the following manner from a reference
state, $\left\vert \Psi _{ref}\right\rangle $, that is nonzero $a.e.$ 
\begin{equation}
\left\vert \Psi \left( \nu ,\omega \right) \right\rangle =\breve{T}\left(
\nu ,\omega \right) \left\vert \Psi _{ref}\right\rangle =\breve{\nu}_{N}e^{i%
\breve{\omega}_{N}}\left\vert \Psi _{ref}\right\rangle
\end{equation}%
where 
\begin{eqnarray}
\left\langle \mathsf{z}^{N}|\Psi _{ref}\right\rangle &=&\varrho \left( 
\mathsf{z}^{N}\right) e^{iS\left( \mathsf{z}^{N}\right) }  \notag \\
\breve{\nu}_{N} &=&\int \nu \left( \mathsf{z}^{N}\right) \left\vert \mathsf{z%
}^{N}\right\rangle \left\langle \mathsf{z}^{N}\right\vert d\mathsf{z}^{N}\in 
\mathcal{B}\left( \mathcal{H}^{N}\right) \\
\breve{\omega}_{N} &=&\int \omega \left( \mathsf{z}^{N}\right) \left\vert 
\mathsf{z}^{N}\right\rangle \left\langle \mathsf{z}^{N}\right\vert d\mathsf{z%
}^{N}\in \mathcal{B}\left( \mathcal{H}^{N}\right)
\end{eqnarray}%
and 
\begin{equation}
\left\vert \mathsf{z}^{N}\right\rangle =\left\vert \mathsf{z}_{1}\cdots 
\mathsf{z}_{N}\right\rangle ;\quad \mathsf{z}_{i}\neq \mathsf{z}_{j},\quad
1\leq i<j\leq N
\end{equation}%
$\lceil $In this formulation it does not matter if the coordinates are equal
or not as for any $N$-particle operator $\breve{\varsigma}_{N}$ 
\begin{equation}
\breve{\varsigma}_{N}=\int \varsigma \left( \mathsf{z}^{N}\right)
\prod_{1\leq i<j\leq N}\delta ^{C}\left( \mathsf{z}_{i}-\mathsf{z}%
_{j}\right) \left\vert \mathsf{z}^{N}\right\rangle \left\langle \mathsf{z}%
^{N}\right\vert d\mathsf{z}^{N}=\int \varsigma \left( \mathsf{z}^{N}\right)
\left\vert \mathsf{z}^{N}\right\rangle \left\langle \mathsf{z}%
^{N}\right\vert d\mathsf{z}^{N}
\end{equation}%
where $\delta ^{C}\left( \mathsf{z}_{i}-\mathsf{z}_{j}\right) =1-\delta
\left( \mathsf{z}_{i}-\mathsf{z}_{j}\right) $, as $\left\vert \mathsf{z}%
_{1}\cdots \mathsf{z}_{N}\right\rangle =0$ if $\mathsf{z}_{i}\neq \mathsf{z}%
_{j}$ for any $i\neq j$.$\rfloor $ Then clearly 
\begin{equation}
I_{N}=\int \left\vert \Psi \left( \nu ,\omega \right) \right\rangle
\left\langle \Psi \left( \nu ,\omega \right) \right\vert d\mu \left( \nu
,\omega \right) =\int \breve{T}\left( \nu ,\omega \right) \left\vert \Psi
_{ref}\right\rangle \left\langle \Psi _{ref}\right\vert \breve{T}\left( \nu
,\omega \right) ^{^{\dagger }}d\mu \left( \nu ,\omega \right)
\end{equation}%
where 
\begin{equation}
d\mu \left( \nu ,\omega \right) =\frac{1}{\left\langle \Psi \left( \nu
,\omega \right) |\Psi \left( \nu ,\omega \right) \right\rangle }d\mu
_{r}\left( \nu \right) d\mu _{i}\left( \omega \right)
\end{equation}%
\qquad \qquad

The strong closure of the space of diagonal integral operators that map $%
\mathcal{H}^{N}\rightarrow \mathcal{H}^{N}$ is the space of integral
operators with symmetric integral kernels $T\left( \mathsf{z}^{N}\right) $ $%
\in L_{\infty }\left( \mathsf{Z}^{N}\right) $ /Walter Thirring, \textit{%
Quantum Mechanics of Atoms and Molecules -A course in Mathematical Physics
3, pp 50-51; }Springer Verlag , New York 1981/. The space $L_{\infty }\left( 
\mathsf{Z}^{N}\right) $ is a $w^{\ast }$-algebra. The following extract from
Thirring should be noted:

\begin{quotation}
$L_{\infty }\left( \mathsf{Z}^{N}\right) ,$\textit{considered as
multiplication operators on} $L_{2}\left( \mathsf{Z}^{N}\right) $ \textit{is
maximally abelian. }\emph{Every function in }$L_{2}\left( \mathsf{Z}%
^{N}\right) $\emph{\ that is nonzero }$a.e.$\emph{\ is a cyclic vector}%
\textit{. Functions that vanish on some set} $\left[ \Omega ^{N}\right] ^{C}=%
\mathsf{Z}^{N}-\Omega ^{N}\subset \mathsf{Z}^{N}$ \textit{form invariant
subspaces. Thus} $L_{\infty }\left( \mathsf{Z}^{N}\right) $ \textit{is
reducible and not a factor.}\emph{\ }
\end{quotation}

One can see that $T\left( \nu ,\omega \right) $ (the integral kernel of $%
\breve{T}\left( \nu ,\omega \right) )$ is a parameterization of a
representation of an abelian $w^{\ast }$-algebra, $\mathfrak{A}$, on $%
L_{2}\left( \mathsf{Z}^{N}\right) $ by $L_{\infty }\left( \mathsf{Z}%
^{N}\right) $. This representation is reducible as $L_{2}\left( \mathsf{Z}%
^{N}\right) $ contains the following invariant subspaces 
\begin{equation}
L_{2}\left( \Omega ^{N}\right) =\left\{ \Psi |\Psi \left( \mathsf{z}%
^{N}\right) =0\,\text{if \thinspace }\mathsf{z}^{N}\notin \Omega ^{N}\,\text{%
,\thinspace }\int_{\Omega ^{N}}d\mathsf{z}^{N}\neq 0\,\text{and\thinspace }%
\Omega ^{N}\subset \mathsf{Z}^{N}\right\}
\end{equation}%
The operator $\breve{T}\left( \nu ,\omega \right) $ produces elements that
have amplitudes that belong to $L_{2}\left( \Omega ^{N}\right) $ when Support%
$\left( \nu \right) =$ $\Omega ^{N}$.

The set of invertible diagonal operators in $\mathfrak{A}$ form a group, $%
\mathfrak{G}_{pos}$, that also has a reducible representation on $L_{\infty
}\left( \mathsf{Z}^{N}\right) $, as clearly $L_{2}\left( \Omega ^{N}\right) $
are also invariant subspaces of $\mathfrak{G}_{pos}$. However \emph{it is
not true} that every function in $L_{2}\left( \mathsf{Z}^{N}\right) $\ that
is nonzero $a.e.$\ is a cyclic vector for $\mathfrak{G}_{pos}$.

$\lceil $This observation also holds in the discrete case, consider the
abelian algebra generated by $\left\{ \tsum\limits_{1\leq i\leq \infty
}\lambda _{i}\left\vert e_{i}\right\rangle \left\langle e_{i}\right\vert
\right\} $ acting on a vector $a=\tsum\limits_{1\leq i\leq \infty }\alpha
_{i}\left\vert e_{i}\right\rangle $ that such that $\alpha _{i}\neq 0\forall
i\rfloor $

From the preceding we have that $L_{\infty }\left( \mathsf{Z}^{N}\right) $
acting on an $a.e.$ nonzero amplitude produces all (unnormalized) vector
states, the question thus arrises on whether $L_{\infty }\left( \mathsf{Z}%
^{N}\right) $ can produce all unnormalized states by acting on an a $a.e.$%
nonzero positive kernel $D_{ref}\left( \mathsf{z}^{N},\mathsf{z}^{\prime
N}\right) $ of an operator $\in \mathcal{B}_{1+}\left( \mathcal{H}%
^{N}\right) $ in the following manner%
\begin{equation}
D\left( \mathsf{z}^{N},\mathsf{z}^{\prime N}\right) =T\left( \mathsf{z}%
^{N}\right) D_{ref}\left( \mathsf{z}^{N},\mathsf{z}^{\prime N}\right)
T\left( \mathsf{z}^{\prime N}\right) ^{\ast }
\end{equation}%
One can start by examining the diagonal elements to obtain 
\begin{eqnarray}
D\left( \mathsf{z}^{N},\mathsf{z}^{N}\right) &=&\left\vert T\left( \mathsf{z}%
^{N}\right) \right\vert ^{2}D_{ref}\left( \mathsf{z}^{N},\mathsf{z}%
^{N}\right)  \notag \\
&\Rightarrow &\left\vert T\left( \mathsf{z}^{N}\right) \right\vert =\left\{ 
\begin{array}{c}
\left( \frac{D\left( \mathsf{z}^{N},\mathsf{z}^{N}\right) }{D_{ref}\left( 
\mathsf{z}^{N},\mathsf{z}^{N}\right) }\right) ^{\frac{1}{2}};\quad
D_{ref}\left( \mathsf{z}^{N},\mathsf{z}^{N}\right) \neq 0 \\ 
0;\qquad \qquad \qquad \quad \quad D_{ref}\left( \mathsf{z}^{N},\mathsf{z}%
^{N}\right) =0\text{ }%
\end{array}%
\right. \text{ }
\end{eqnarray}%
which is well defined $a.e.$ The off diagonal elements will "determine" the
phase of the transformation $T\left( \mathsf{z}^{N}\right) $ $=\left\vert
T\left( \mathsf{z}^{N}\right) \right\vert e^{i\theta \left( \mathsf{z}%
^{N}\right) }$ by 
\begin{eqnarray}
D\left( \mathsf{z}^{N},\mathsf{z}^{\prime N}\right) &=&\left\vert T\left( 
\mathsf{z}^{N}\right) T\left( \mathsf{z}^{\prime N}\right) \right\vert
e^{i\left( \theta \left( \mathsf{z}^{N}\right) -\theta \left( \mathsf{z}%
^{\prime N}\right) \right) }D_{ref}\left( \mathsf{z}^{N},\mathsf{z}^{\prime
N}\right)  \notag \\
&\Rightarrow &e^{i\left( \theta \left( \mathsf{z}^{N}\right) -\theta \left( 
\mathsf{z}^{\prime N}\right) \right) }=\frac{D\left( \mathsf{z}^{N},\mathsf{z%
}^{\prime N}\right) }{D_{ref}\left( \mathsf{z}^{N},\mathsf{z}^{\prime
N}\right) \left\vert T\left( \mathsf{z}^{N}\right) T\left( \mathsf{z}%
^{\prime N}\right) \right\vert }
\end{eqnarray}%
which is well defined as long as 
\begin{equation}
\left\vert T\left( \mathsf{z}^{N}\right) T\left( \mathsf{z}^{\prime
N}\right) \right\vert =0\Longleftrightarrow D\left( \mathsf{z}^{N},\mathsf{z}%
^{\prime N}\right) =0
\end{equation}%
which means that 
\begin{equation}
D\left( \mathsf{z}^{\prime N},\mathsf{z}^{\prime N}\right) =0\text{ or }%
D\left( \mathsf{z}^{N},\mathsf{z}^{N}\right) =0\Longrightarrow D\left( 
\mathsf{z}^{N},\mathsf{z}^{\prime N}\right) =0
\end{equation}%
when $D\left( \mathsf{z}^{N},\mathsf{z}^{\prime N}\right) \geq 0$. First we
see if is true for proper vectors, in that case 
\begin{eqnarray}
D\left( \mathsf{v}^{N},\mathsf{v}^{N}\right) &=&\left\langle \mathsf{v}^{N}|%
\breve{D}\mathsf{v}^{N}\right\rangle =\left\langle \mathsf{v}^{N}|\breve{Q}%
^{\dagger }\breve{Q}\mathsf{v}^{N}\right\rangle =0  \notag \\
&\Rightarrow &\breve{Q}\left\vert \mathsf{v}^{N}\right\rangle =0  \label{c1}
\end{eqnarray}%
and 
\begin{eqnarray}
D\left( \mathsf{v}^{\prime N},\mathsf{v}^{\prime N}\right) &=&\left\langle 
\mathsf{v}^{\prime N}|\breve{D}\mathsf{v}^{\prime N}\right\rangle
=\left\langle \mathsf{v}^{\prime N}|\breve{Q}^{\dagger }\breve{Q}\mathsf{v}%
^{\prime N}\right\rangle =0  \notag \\
&\Rightarrow &\breve{Q}\left\vert \mathsf{v}^{\prime N}\right\rangle =0
\label{c2}
\end{eqnarray}%
so either condition eq. $\left( \ref{c1}\right) $ or eq. $\left( \ref{c2}%
\right) $ clearly implies that 
\begin{equation}
\left\langle \mathsf{v}^{N}|\breve{Q}^{\dagger }\breve{Q}\mathsf{v}^{\prime
N}\right\rangle =D\left( \mathsf{v}^{N},\mathsf{v}^{\prime N}\right) =0
\end{equation}%
To see if this true for the generalized basis we must see if 
\begin{equation}
\left\langle \mathsf{z}^{N}|\breve{Q}^{\dagger }\breve{Q}\mathsf{z}%
^{N}\right\rangle =0\Rightarrow \breve{Q}\left\vert \mathsf{z}%
^{N}\right\rangle =0
\end{equation}%
is true in a rigged Hilbert space. $\lceil \mathcal{V}_{L}\supset \mathcal{%
H\supset }\mathcal{V}_{S};$ $\mathcal{V}_{L}$ and $\mathcal{V}_{S}$ are
Banach spaces and $\mathcal{V}_{S}$ is dense in $\mathcal{H}$; $\breve{Q}$
acting on $\mathcal{V}_{L}$ is defined by 
\begin{equation}
\left( \breve{Q}^{\ddagger }\mathsf{z}|\varphi \right) =\left( \mathsf{z}|%
\breve{Q}\varphi \right) \forall \left\vert \varphi \right\rangle \in 
\mathcal{V}_{S}
\end{equation}%
and $\breve{Q}^{\ddagger }$ is an unique extension of $\breve{Q}$. Note that 
$\breve{Q}\left\vert \mathsf{z}^{N}\right\rangle =0$ means 
\begin{equation}
\left( \mathsf{z}|\breve{Q}\varphi \right) =0\forall \left\vert \varphi
\right\rangle \in \mathcal{V}_{S}
\end{equation}%
which implies that $\breve{Q}^{\ddagger }\left\vert \mathsf{z}%
^{N}\right\rangle $ is the zero vector in $\mathcal{V}_{L}$.$\rfloor $ Thus $%
\left\langle \mathsf{z}^{N}|\breve{Q}^{\dagger }\breve{Q}\mathsf{z}%
^{N}\right\rangle =0$ or $\left\langle \mathsf{z}^{\prime N}|\breve{Q}%
^{\dagger }\breve{Q}\mathsf{z}^{\prime N}\right\rangle =0\Rightarrow
\left\langle \mathsf{z}^{\prime N}|\breve{Q}^{\dagger }\breve{Q}\mathsf{z}%
^{N}\right\rangle =0$. Thus we have that 
\begin{equation}
\left\vert T\left( \mathsf{z}^{N}\right) T\left( \mathsf{z}^{\prime
N}\right) \right\vert =0\Longleftrightarrow D\left( \mathsf{z}^{N},\mathsf{z}%
^{\prime N}\right) =0
\end{equation}%
and we can define phase differences as 
\begin{equation}
e^{i\left( \theta \left( \mathsf{z}^{N}\right) -\theta \left( \mathsf{z}%
^{\prime N}\right) \right) }=\left\{ 
\begin{array}{c}
\frac{D\left( \mathsf{z}^{N},\mathsf{z}^{\prime N}\right) }{D_{ref}\left( 
\mathsf{z}^{N},\mathsf{z}^{\prime N}\right) \left\vert T\left( \mathsf{z}%
^{N}\right) T\left( \mathsf{z}^{\prime N}\right) \right\vert };\left\vert
T\left( \mathsf{z}^{N}\right) T\left( \mathsf{z}^{\prime N}\right)
\right\vert \neq 0,D_{ref}\left( \mathsf{z}^{N},\mathsf{z}^{\prime N}\right)
\neq 0 \\ 
\text{\qquad \qquad }0\qquad \text{\qquad \qquad \qquad otherwise\qquad
\qquad \qquad \qquad \qquad }%
\end{array}%
\right.
\end{equation}%
which is sufficient to define the transformation $D_{ref}\rightarrow D$ 
\emph{for any }$D$ as long as $D_{ref}\neq 0\,a.e.$

The transformations $\hat{T}$ can either decrease or leave invariant the
rank of $\breve{D}^{N}$, thus only if $\breve{D}_{ref}$ has maximal rank
will all density operators be generated by such transformations ($D_{ref}$
need not be a density operator kernel associated with a fixed number of
particles, all that is needed is that $D_{ref}$ $\in \mathcal{B}_{1+}\left( 
\mathcal{H}^{N}\right) $ with $rank\left( D_{ref}\right) =\infty $ and thus
could be the kernel of any density operator, including ones describing
indeterminate number of particles.)

\subsection{Second Quantized form of Diagonal Operators}

The operator 
\begin{equation}
\breve{T}_{N}=\int T\left( \mathsf{z}^{N}\right) \left\vert \mathsf{z}%
^{N}\right\rangle \left\langle \mathsf{z}^{N}\right\vert d\mathsf{z}^{N}
\end{equation}%
is the restriction of the second quantized operator%
\begin{equation}
\breve{T}=\int T\left( \mathsf{z}^{N}\right) \Phi \left( \mathsf{z}%
^{N}\right) ^{\dagger }\Phi \left( \mathsf{z}^{N}\right) d\mathsf{z}^{N}
\end{equation}%
to $\mathcal{H}^{N}$, where 
\begin{equation}
\Phi \left( \mathsf{z}^{N}\right) =\Phi \left( \mathsf{z}_{N}\right) \cdots
\Phi \left( \mathsf{z}_{1}\right)
\end{equation}%
and 
\begin{equation}
\Phi \left( \mathsf{z}\right) ^{\dagger }\left\vert \phi \right\rangle
=\left\vert \mathsf{z}\right\rangle
\end{equation}%
As%
\begin{eqnarray}
\Phi \left( \mathsf{z}^{N}\right) ^{\dagger }\Phi \left( \mathsf{z}%
^{N}\right) &=&\Phi \left( \mathsf{z}_{1}\right) ^{\dagger }\cdots \Phi
\left( \mathsf{z}_{N}\right) ^{\dagger }\Phi \left( \mathsf{z}_{N}\right)
\cdots \Phi \left( \mathsf{z}_{1}\right)  \notag \\
&=&\prod_{1\leq i<j\leq N}\delta ^{C}\left( \mathsf{z}_{i}-\mathsf{z}%
_{j}\right) \Phi \left( \mathsf{z}_{1}\right) ^{\dagger }\Phi \left( \mathsf{%
z}_{1}\right) \cdots \Phi \left( \mathsf{z}_{N}\right) ^{\dagger }\Phi
\left( \mathsf{z}_{N}\right)
\end{eqnarray}%
we have that%
\begin{eqnarray}
\breve{T} &=&\int T\left( \mathsf{z}^{N}\right) \prod_{1\leq i<j\leq
N}\delta ^{C}\left( \mathsf{z}_{i}-\mathsf{z}_{j}\right) \Phi \left( \mathsf{%
z}_{1}\right) ^{\dagger }\Phi \left( \mathsf{z}_{1}\right) \cdots \Phi
\left( \mathsf{z}_{N}\right) ^{\dagger }\Phi \left( \mathsf{z}_{N}\right) d%
\mathsf{z}^{N}  \notag \\
&=&\int \sum_{k=0}^{N}\left\{ T_{k}\left( \mathsf{z}^{N}\right) \right\}
\prod_{1\leq i<j\leq N}\delta ^{C}\left( \mathsf{z}_{i}-\mathsf{z}%
_{j}\right) \Phi \left( \mathsf{z}_{1}\right) ^{\dagger }\Phi \left( \mathsf{%
z}_{1}\right) \cdots \Phi \left( \mathsf{z}_{N}\right) ^{\dagger }\Phi
\left( \mathsf{z}_{N}\right) d\mathsf{z}^{N}
\end{eqnarray}%
where 
\begin{equation}
T_{k}\left( \mathsf{z}^{N}\right) =\sum_{1\leq i_{1}<i_{2}\cdots <i_{k}\leq
N}T_{k}\left( \mathsf{z}_{i_{1}}\mathsf{z}_{i_{2}}\cdots \mathsf{z}%
_{i_{k}}\right)
\end{equation}%
and all the functions are symmetric.

\section{Outline of Invariance Groups and Subgroups of Kernels of
Contraction Maps}

\subsection{General Description}

A locally compact abelian semi group $\mathfrak{S}$ that sends normalized
states into normalized states can be defined as 
\begin{equation}
\mathfrak{S}=\left\{ \mathcal{T}|\mathcal{T}\left( \breve{X}\right) =\frac{%
\breve{T}\breve{X}\breve{T}^{\dagger }}{Tr\left\{ \breve{T}\breve{X}\breve{T}%
^{\dagger }\right\} };\quad \in \breve{X},\breve{T}\in \mathcal{B}\left( 
\mathcal{H}^{N}\right) \right\}  \label{gd}
\end{equation}%
where 
\begin{equation}
\breve{T}=\breve{\nu}_{N}e^{\breve{\omega}_{N}}=\int \nu \left( \mathsf{z}%
^{N}\right) e^{i\omega \left( \mathsf{z}^{N}\right) }\left\vert \mathsf{z}%
^{N}\right\rangle \left\langle \mathsf{z}^{N}\right\vert d\mathsf{z}%
^{N}=\int T\left( \mathsf{z}^{N}\right) \left\vert \mathsf{z}%
^{N}\right\rangle \left\langle \mathsf{z}^{N}\right\vert d\mathsf{z}^{N}
\end{equation}%
The overcomplete complete basis of generalized vectors $\left\{ \left\vert 
\mathsf{z}^{N}\right\rangle \right\} $ $\equiv \left\{ \left\vert \mathsf{z}%
_{1}\cdots \mathsf{z}_{N}\right\rangle \right\} $ is arbitrary and different
bases give different isomorphic representations of $\mathfrak{S}$. Note that
eq. $\left( \ref{gd}\right) $ \emph{is not linear }i.e. the representation
is not a linear representation of $\mathfrak{S}$. This sub semigroup maps
the convex set of normalized states, $\mathcal{S}_{N}$, to itself, however
it is not defined on the kernel of the contraction maps, which are defined
later.

One can define another semigroup $\mathfrak{S}_{pos}$, which is also an
abelian $w^{\ast }$-algebra, by 
\begin{equation}
\mathfrak{S}_{pos}=\left\{ \hat{T}|\hat{T}\left( \breve{X}\right) =\breve{T}%
\breve{X}\breve{T}^{\dagger };\quad \breve{T}\in \mathcal{B}\left( \mathcal{H%
}^{N}\right) \right\}  \label{pos}
\end{equation}%
This maps the cone of positive operators $\mathcal{B}\left( \mathcal{H}%
^{N}\right) _{+}$ to itself, it also maps the cone of positive trace class
operators to itself. The maps defined in eq. $\left( \ref{pos}\right) $ are
linear maps, but they do \emph{not }map $\mathcal{S}_{N}$ to itself as 
\begin{equation}
Tr\left\{ \breve{T}\breve{D}^{N}\breve{T}^{\dagger }\right\} \neq 1\text{ }%
\forall \breve{D}^{N}\text{ unless }\breve{T}^{\dagger }\breve{T}=I
\end{equation}%
This semigroup is defined on the kernels of the contraction maps, but does
not leave them invariant, it is the sub semigroups (which are also sub $%
w^{\ast }$-algebras) $\mathfrak{S}_{\Xi }$, $\mathfrak{S}_{L_{N}^{1}}$ and $%
\mathfrak{S}_{L_{N}^{2}}$ of $\mathfrak{S}_{pos}$ that have this property.
Clearly 
\begin{equation}
\mathfrak{S}_{\Xi }\supset \mathfrak{S}_{L_{N}^{1}}\supset \mathfrak{S}%
_{L_{N}^{2}}
\end{equation}

\begin{lemma}
\label{lemma1}$\breve{X}_{1}=\breve{X}_{2}$ when $\breve{X}_{1},\breve{X}%
_{2}\in \mathcal{B}\left( \mathcal{H}^{N}\right) $ iff $X_{1}\left( \mathsf{z%
}^{N},\mathsf{z}^{\prime N}\right) \overset{a.e.}{=}X_{2}\left( \mathsf{z}%
^{N},\mathsf{z}^{\prime N}\right) $

\begin{proof}
See "Bounded Integral Operators on $L^{2}$ Spaces" by Halmos and Sunder.
\end{proof}
\end{lemma}

\begin{lemma}
Let $\breve{T}_{1}$ and $\breve{T}_{2}\in \mathfrak{S}_{pos}$ then $\breve{T}%
_{1}=\breve{T}_{2}$ iff $T_{1}\left( \mathsf{z}^{N}\right) =T_{2}\left( 
\mathsf{z}^{N}\right) $ a.e. with respect to the measure $d\mathsf{z}^{N}$
\end{lemma}

\begin{proof}
One can use \ref{lemma1} or consider the operator equality is defined by 
\begin{equation}
\left\Vert \left( \breve{T}_{1}-\breve{T}_{2}\right) \right\Vert
=\sup_{\left\Vert v\right\Vert =1}\left\Vert \left( \breve{T}_{1}-\breve{T}%
_{2}\right) v\right\Vert =0;\qquad v\in \mathcal{H}^{N}
\end{equation}%
i.e. 
\begin{equation}
\int \left\vert T_{1}\left( \mathsf{z}^{N}\right) -T_{2}\left( \mathsf{z}%
^{N}\right) \right\vert ^{2}\left\vert v\left( \mathsf{z}^{N}\right)
\right\vert ^{2}d\mathsf{z}^{N}=0\qquad \forall v\in \mathcal{H}^{N}
\end{equation}%
this implies that 
\begin{equation}
\left\vert T_{1}\left( \mathsf{z}^{N}\right) -T_{2}\left( \mathsf{z}%
^{N}\right) \right\vert ^{2}\overset{a.e.}{=}0
\end{equation}%
which implies that 
\begin{equation}
T_{1}\left( \mathsf{z}^{N}\right) -T_{2}\left( \mathsf{z}^{N}\right) \overset%
{a.e.}{=}0
\end{equation}
\end{proof}

This Lemma allows us to uniquely define the diagonal operator component, $%
\breve{X}_{d}$ , of the trace class operator $\breve{X}$ as 
\begin{equation}
\breve{X}_{d}=\int X\left( \mathsf{z}^{N}\right) \left\vert \mathsf{z}%
^{N}\right\rangle \left\langle \mathsf{z}^{N}\right\vert d\mathsf{z}^{N}
\end{equation}%
where 
\begin{eqnarray}
X\left( \mathsf{z}^{N}\right) &=&X\left( \mathsf{z}^{N},\mathsf{z}^{N}\right)
\notag \\
\breve{X} &=&\int X\left( \mathsf{z}^{N},\mathsf{z}^{\prime N}\right) d%
\mathsf{z}^{N}d\mathsf{z}^{\prime N}
\end{eqnarray}%
and we have identified operator kernels that are equal almost everywhere.
Clearly $\breve{X}_{d}\in \mathcal{B}_{1}\left( \mathcal{H}^{N}\right) $ if $%
\breve{X}$ $\in \mathcal{B}_{1}\left( \mathcal{H}^{N}\right) $ as 
\begin{equation}
Tr\left\{ \left\{ \breve{X}^{\dagger }\breve{X}\right\} ^{\frac{1}{2}%
}\right\} =Tr\left\{ \left\{ U\breve{x}^{\dagger }\breve{x}U^{\dagger
}\right\} ^{\frac{1}{2}}\right\} =Tr\left\{ \left\{ \breve{x}^{\dagger }%
\breve{x}\right\} ^{\frac{1}{2}}\right\} =\int \left\vert x\left( \mathsf{%
\eta }^{N}\right) \right\vert d\mathsf{\eta }^{N}
\end{equation}%
and 
\begin{equation}
Tr\left\{ \left\{ \breve{X}_{d}^{\dagger }\breve{X}_{d}\right\} ^{\frac{1}{2}%
}\right\} =\int \left\vert X_{d}\left( \mathsf{z}^{N}\right) \right\vert d%
\mathsf{z}^{N}
\end{equation}%
Which gives that 
\begin{equation}
Tr\left\{ \left\{ \breve{X}^{\dagger }\breve{X}\right\} ^{\frac{1}{2}%
}\right\} \geq Tr\left\{ \left\{ \breve{X}_{d}^{\dagger }\breve{X}%
_{d}\right\} ^{\frac{1}{2}}\right\}
\end{equation}%
Thus $\breve{X}_{od}\in \mathcal{B}_{1}\left( \mathcal{H}^{N}\right) $ and
is uniquely defined by 
\begin{equation}
\breve{X}_{od}=\breve{X}-\breve{X}_{d}
\end{equation}%
as $\mathcal{B}_{1}\left( \mathcal{H}^{N}\right) $ is a linear space and 
\begin{equation}
\breve{X}_{od}=\int X_{od}\left( \mathsf{z}^{N},\mathsf{z}^{\prime N}\right)
\left\vert \mathsf{z}^{N}\right\rangle \left\langle \mathsf{z}^{\prime
N}\right\vert \delta ^{C}\left( \mathsf{z}^{N}-\mathsf{z}^{\prime N}\right) d%
\mathsf{z}^{N}d\mathsf{z}^{\prime N}
\end{equation}

Note that $\mathcal{B}_{1}\left( \mathcal{H}^{N}\right) \subseteq \mathcal{B}%
_{2}\left( \mathcal{H}^{N}\right) $ and in fact $\mathcal{B}_{2}\left( 
\mathcal{H}^{N}\right) =\mathcal{B}_{1}\left( \mathcal{H}^{N}\right) ^{2}$,
it is clear that 
\begin{equation}
Tr\left\{ \breve{X}^{\dagger }\breve{X}\right\} \geq Tr\left\{ \breve{X}%
_{d}^{\dagger }\breve{X}_{d}\right\}
\end{equation}%
so that $\breve{X}\in \mathcal{B}_{2}\left( \mathcal{H}^{N}\right)
\Rightarrow \breve{X}_{d}\in \mathcal{B}_{2}\left( \mathcal{H}^{N}\right) $.
Further $\breve{X}_{d}\bot \breve{Y}_{od}$ $\forall \breve{Y}_{od}\in 
\mathcal{B}_{2}\left( \mathcal{H}^{N}\right) $, thus it is also true for
those parts belonging to $\mathcal{B}_{1}\left( \mathcal{H}^{N}\right) $.

\begin{theorem}
If $\breve{T}_{1}\breve{X}\breve{T}_{1}^{\dagger }=\breve{T}_{2}\breve{X}%
\breve{T}_{2}^{\dagger }$ then $T_{1}\left( \mathsf{z}^{N}\right)
T_{1}\left( \mathsf{z}^{\prime N}\right) ^{\ast }\overset{a.e.\left(
X\right) }{=}T_{2}\left( \mathsf{z}^{N}\right) T_{2}\left( \mathsf{z}%
^{\prime N}\right) ^{\ast }$

\begin{proof}
$T_{1}\left( \mathsf{z}^{N}\right) T_{1}\left( \mathsf{z}^{\prime N}\right)
^{\ast }X\left( \mathsf{z}^{N},\mathsf{z}^{\prime N}\right) d\mathsf{z}^{N}d%
\mathsf{z}^{\prime N}\overset{a.e.}{=}T_{2}\left( \mathsf{z}^{N}\right)
T_{2}\left( \mathsf{z}^{\prime N}\right) ^{\ast }X\left( \mathsf{z}^{N},%
\mathsf{z}^{\prime N}\right) d\mathsf{z}^{N}d\mathsf{z}^{\prime N}$
\end{proof}
\end{theorem}

\begin{corollary}
$\breve{T}_{1}\breve{X}\breve{T}_{1}^{\dagger }\neq \breve{T}_{2}\breve{X}%
\breve{T}_{2}^{\dagger }$ when $\breve{T}_{1}\neq \breve{T}_{2}$ unless $%
T_{1}\left( \mathsf{z}^{N}\right) \overset{a.e.\left( X\right) }{=}%
T_{2}\left( \mathsf{z}^{N}\right) $
\end{corollary}

One can define the following linear contraction maps

\begin{description}
\item 
\begin{itemize}
\item The density contraction%
\begin{eqnarray}
\Xi &:&\mathcal{B}_{1}\left( \mathcal{H}^{N}\right) \rightarrow L_{1}\left( 
\mathsf{Z}_{1}\right)  \notag \\
\Xi \left( \breve{X}\right) \left( \mathsf{z}_{1}\right) &=&Tr\left\{ \breve{%
X}\Phi \left( \mathsf{z}_{1}\right) ^{\dagger }\Phi \left( \mathsf{z}%
_{1}\right) \right\}  \label{d0}
\end{eqnarray}%
giving 
\begin{eqnarray}
\Xi \left( \breve{X}\right) &=&\int \Xi \left( \breve{X}\right) \left( 
\mathsf{z}_{1}\right) \left\vert \mathsf{z}_{1}\right\rangle \left\langle 
\mathsf{z}_{1}\right\vert d\mathsf{z}_{1}  \notag \\
&=&\int X\left( \mathsf{z}_{1}\mathsf{z}_{2}\cdots \mathsf{z}_{N},\mathsf{z}%
_{1}\mathsf{z}_{2}\cdots \mathsf{z}_{N}\right) \left\vert \mathsf{z}%
_{1}\right\rangle \left\langle \mathsf{z}_{1}^{\prime }\right\vert d\mathsf{z%
}_{1}d\mathsf{z}_{2}\cdots d\mathsf{z}_{N}
\end{eqnarray}%
where we have used embedding of $L_{1}\left( \mathsf{Z}_{1}\right) $ in $%
\mathcal{B}_{1}\left( \mathcal{H}^{1}\right) $. The map $\Xi $ \emph{basis
dependent }so should really be denoted as $\Xi _{\mathsf{Z}}$. If $\mathsf{%
Z\neq Z}^{\prime }$ then $\Xi _{\mathsf{Z}}\left( \breve{X}\right) \neq \Xi
_{\mathsf{Z}^{\prime }}\left( \breve{X}\right) $ and 
\begin{equation}
\Xi _{\mathsf{Z}}\left( U_{N}\breve{X}U_{N}^{\dagger }\right) \neq U_{1}\Xi
_{\mathsf{Z}}\left( \breve{X}\right) U_{1}^{\dagger }
\end{equation}

\item The one density contraction 
\begin{eqnarray}
&&L_{N}^{1}:\mathcal{B}_{1}\left( \mathcal{H}^{N}\right) \rightarrow 
\mathcal{B}_{1}\left( \mathcal{H}^{1}\right)  \notag \\
&&\Xi \left( \breve{X}\right) \left( \mathsf{z}_{1},\mathsf{z}_{1}^{\prime
}\right) =Tr\left\{ \breve{X}\Phi \left( \mathsf{z}_{1}^{\prime }\right)
^{\dagger }\Phi \left( \mathsf{z}_{1}\right) \right\}
\end{eqnarray}%
giving 
\begin{eqnarray}
L_{N}^{1}\left( \breve{X}\right) &=&\int \Xi \left( \breve{X}\right) \left( 
\mathsf{z}_{1},\mathsf{z}_{1}^{\prime }\right) \left\vert \mathsf{z}%
_{1}\right\rangle \left\langle \mathsf{z}_{1}^{\prime }\right\vert d\mathsf{z%
}_{1}d\mathsf{z}_{1}^{\prime }  \notag \\
\left. {}\right. \qquad &=&\int X\left( \mathsf{z}_{1}\mathsf{z}_{2}\cdots 
\mathsf{z}_{N},\mathsf{z}_{1}^{\prime }\mathsf{z}_{2}\cdots \mathsf{z}%
_{N}\right) \left\vert \mathsf{z}_{1}\right\rangle \left\langle \mathsf{z}%
_{1}^{\prime }\right\vert d\mathsf{z}_{1}d\mathsf{z}_{1}^{\prime }d\mathsf{z}%
_{2}\cdots d\mathsf{z}_{N}
\end{eqnarray}

\item The two density contraction%
\begin{eqnarray}
L_{N}^{2} &:&\mathcal{B}_{1}\left( \mathcal{H}^{N}\right) \rightarrow 
\mathcal{B}_{1}\left( \mathcal{H}^{2}\right)  \notag \\
\Xi \left( \breve{X}\right) \left( \mathsf{z}_{1},\mathsf{z}_{2},\mathsf{z}%
_{1}^{\prime },\mathsf{z}_{2}^{\prime }\right) &=&Tr\left\{ \breve{X}\Phi
\left( \mathsf{z}_{1}^{\prime }\right) ^{\dagger }\Phi \left( \mathsf{z}%
_{2}^{\prime }\right) ^{\dagger }\Phi \left( \mathsf{z}_{2}\right) \Phi
\left( \mathsf{z}_{1}\right) \right\}  \label{d2}
\end{eqnarray}%
giving 
\begin{eqnarray}
&&L_{N}^{2}\left( \breve{X}\right) =\int \Xi \left( \breve{X}\right) \left( 
\mathsf{z}_{1},\mathsf{z}_{2},\mathsf{z}_{1}^{\prime },\mathsf{z}%
_{2}^{\prime }\right) \left\vert \mathsf{z}_{1}\mathsf{z}_{2}\right\rangle
\left\langle \mathsf{z}_{1}^{\prime }\mathsf{z}_{2}^{\prime }\right\vert d%
\mathsf{z}_{1}d\mathsf{z}_{2}d\mathsf{z}_{1}^{\prime }d\mathsf{z}%
_{2}^{\prime }  \notag \\
&=&\int X\left( \mathsf{z}_{1}\mathsf{z}_{2}\mathsf{z}_{3}\cdots \mathsf{z}%
_{N},\mathsf{z}_{1}^{\prime }\mathsf{z}_{2}^{\prime }\mathsf{z}_{3}\cdots 
\mathsf{z}_{N}\right) \left\vert \mathsf{z}_{1}\mathsf{z}_{2}\right\rangle
\left\langle \mathsf{z}_{1}^{\prime }\mathsf{z}_{2}^{\prime }\right\vert d%
\mathsf{z}_{1}d\mathsf{z}_{2}d\mathsf{z}_{1}^{\prime }d\mathsf{z}%
_{2}^{\prime }d\mathsf{z}_{3}\cdots d\mathsf{z}_{N}
\end{eqnarray}
\end{itemize}

\item \qquad These maps could be defined in the same wrt to any overcomplete
(orthonormal) basis. However only $L_{N}^{1}$ and $L_{N}^{2}$ are bases
independent.
\end{description}

These overcomplete (orthonormal) bases are connected to each other by
unitary transformations 
\begin{equation}
\left\vert \mathsf{z}_{1}\right\rangle =\int U_{1}\left( \mathsf{z}_{1}%
\mathsf{,y}_{1}\right) \left\vert \mathsf{y}_{1}\right\rangle d\mathsf{y}_{1}
\end{equation}%
and we have the intertwining property 
\begin{equation}
L_{N}^{p}\left( U_{N}XU_{N}^{\dagger }\right) =U_{P}L_{N}^{p}\left( X\right)
U_{P}^{\dagger }
\end{equation}%
where 
\begin{equation}
U_{K}=\wedge ^{K}U
\end{equation}%
However as mentioned before 
\begin{equation}
\Xi \left( U_{N}\breve{X}U_{N}^{\dagger }\right) \neq U_{1}\Xi \left( \breve{%
X}\right) U_{1}^{\dagger }
\end{equation}

\qquad We can decompose $\mathcal{B}_{1}\left( \mathcal{H}^{N}\right) $ into
a direct sum of subspaces 
\begin{eqnarray}
\mathcal{B}_{1}\left( \mathcal{H}^{N}\right)  &=&\bigoplus_{i=0}^{N}\mathcal{%
B}_{i}^{N}  \notag \\
&=&Ker\Xi \oplus \left[ Ker\Xi \right] ^{C}  \notag \\
&=&Ker\Xi _{OD}\oplus Ker\Xi _{D}\oplus \left[ Ker\Xi \right] ^{C}
\end{eqnarray}%
where%
\begin{eqnarray}
\mathcal{B}_{i}^{N} &=&KerL_{N}^{i-1}\circleddash KerL_{N}^{i};\qquad 1\leq
i\leq N  \notag \\
\mathcal{B}_{0}^{N} &=&\left[ KerL_{N}^{0}\right] ^{C}
\end{eqnarray}%
and 
\begin{eqnarray}
Ker\Xi _{OD} &=&\left\{ \breve{Y}|\breve{Y}=\int Y\left( \mathsf{z}^{N},%
\mathsf{z}^{N\prime }\right) \delta ^{C}\left( \mathsf{z}^{N}-\mathsf{z}%
^{N\prime }\right) \left\vert \mathsf{z}^{N}\right\rangle \left\langle 
\mathsf{z}^{N\prime }\right\vert d\mathsf{z}^{N}d\mathsf{z}^{N\prime
}\right\}   \notag \\
Ker\Xi _{D} &=&\left\{ \breve{Y}|\breve{Y}=\int Y\left( \mathsf{z}%
^{N}\right) \left\vert \mathsf{z}^{N}\right\rangle \left\langle \mathsf{z}%
^{N}\right\vert d\mathsf{z}^{N};\quad \int Y\left( \mathsf{z}^{N}\right) d%
\mathsf{z}^{N}=0\right\} 
\end{eqnarray}%
Observe that 
\begin{equation}
KerL_{N}^{0}\circleddash Ker\Xi 
\end{equation}%
contains all trace class operators that produce non zero functions that
integrate to zero. The subspaces have the following property 
\begin{equation}
KerL_{N}^{0}\supset Ker\Xi \supset KerL_{N}^{1}\supset KerL_{N}^{2}\supset
\cdots \supset KerL_{N}^{N-1}\supset KerL_{N}^{N}
\end{equation}%
$\lceil $Note that $L_{N}^{N}=I\rfloor .$ The canonical compliment, $Ker\Xi
^{C}$ of $Ker\Xi $ in $\mathcal{B}_{1}\left( \mathcal{H}^{N}\right) $ is
isomorphic to $L_{1}\left( \mathsf{Z}_{1}\right) $ i.e.%
\begin{equation}
\left[ Ker\Xi \right] ^{C}\cong L_{1}\left( \mathsf{Z}_{1}\right) 
\end{equation}%
\begin{equation}
\left[ KerL_{N}^{0}\right] ^{C}\subset \left[ Ker\Xi \right] ^{C}
\end{equation}%
The subspace $\left[ KerL_{N}^{0}\right] ^{C}$ contains trace operators that
have non zero trace, while $\left[ Ker\Xi \right] ^{C}$ can contain
operators that have zero trace. An important theorem for our treatment is

\begin{theorem}
If $\Theta \left( X\right) =\Theta \left( X^{\prime }\right) $ then $%
Tr\left\{ X\right\} =Tr\left\{ X^{\prime }\right\} $

\begin{proof}
The trace of parts of an operator belonging to $Ker\Xi $ are zero and if $%
\Theta \left( X\right) =\Theta \left( X^{\prime }\right) $ the operators
have identical components in $\left[ Ker\Xi \right] ^{C}$ thus their trace
must be identical.

\begin{corollary}
If $\Theta \left( X\right) =\Theta \left( X^{\prime }\right) $ then $\Theta
\left( \frac{X}{Tr\left\{ X\right\} }\right) =\Theta \left( \frac{X^{\prime }%
}{Tr\left\{ X^{\prime }\right\} }\right) $
\end{corollary}
\end{proof}
\end{theorem}

Note that $\Theta \left( \frac{X}{Tr\left\{ X\right\} }\right) =\Theta
\left( \frac{X^{\prime }}{Tr\left\{ X^{\prime }\right\} }\right) $ $%
\not\Rightarrow $ $\Theta \left( X\right) =\Theta \left( X^{\prime }\right) $%
, so if two states contract to the same density it does not mean that
unnormalized states that are scalar multiples of these states also contract
to the same density.

In the finite dimensional case $\left[ KerL_{N}^{0}\right] ^{C}=\left\{
\alpha I|\alpha \in \mathbb{C}\right\} $ as%
\begin{equation}
\mathcal{B}\left( \mathcal{H}^{N}\right) =\mathcal{B}_{0}\left( \mathcal{H}%
^{N}\right) =\mathcal{B}_{1}\left( \mathcal{H}^{N}\right) =\mathcal{B}%
_{2}\left( \mathcal{H}^{N}\right)
\end{equation}%
and 
\begin{equation}
\left[ KerL_{N}^{0}\right] ^{C}=\left[ KerL_{N}^{0}\right] ^{\bot _{HS}}
\end{equation}%
and 
\begin{eqnarray}
\left[ KerL_{N}^{0}\right] ^{\bot _{HS}} &=&\left\{ X|\left( X|Y\right)
_{HS};\forall Y\in KerL_{N}^{0}\equiv \forall Y\backepsilon TrY=0\right\} 
\notag \\
&\Rightarrow &X=\alpha I
\end{eqnarray}%
This together with the fact that for any operator $A$ 
\begin{equation}
A=\left( A-\left( TrA\right) I\right) +\left( TrA\right) I
\end{equation}%
proves the assertion. In the infinite dimensional case $I\notin \mathcal{B}%
_{2}\left( \mathcal{H}^{N}\right) ,$so we have instead 
\begin{equation}
\left[ KerL_{N}^{0}\right] ^{\bot _{HS}}=\left\{ X|\left( X|Y\right)
_{HS};\forall Y\in KerL_{N}^{0}\equiv \forall Y\backepsilon TrY=0\right\}
\end{equation}%
and the question remains whether $\exists X$ with this property and whether
if there $\exists $ any $X$ they belong to $\mathcal{B}_{2}\left( \mathcal{H}%
^{N}\right) $.

\subsection{Explicit Descriptions of the Invariance Subalgebras}

\begin{itemize}
\item The density contraction
\end{itemize}

In order that $Ker\Xi \rightarrow Ker\Xi $ by $\mathfrak{S}_{\Xi }$ we must
have that 
\begin{equation}
\int X\left( \mathsf{z}^{N}\right) \left\vert T\left( \mathsf{z}^{N}\right)
\right\vert ^{2}d\mathsf{z}_{2}\cdots d\mathsf{z}_{N}\overset{a.e.}{=}%
0\qquad \forall \breve{X}\in Ker\Xi
\end{equation}%
i.e. 
\begin{equation}
\hat{T}:Ker\Xi _{D}\rightarrow Ker\Xi _{D}
\end{equation}%
$\lceil $Note that $\hat{T}:Ker\Xi _{OD}\rightarrow Ker\Xi _{OD}$
automatically so this is not a constraint$\rfloor $

\begin{itemize}
\item The one density contraction
\end{itemize}

In order that $KerL_{N}^{1}\rightarrow KerL_{N}^{1}$ by $\mathfrak{S}%
_{L_{N}^{1}}$ we must have that%
\begin{equation}
\int X\left( \mathsf{z}_{1}\mathsf{z}^{N-1}\mathsf{,z}_{1}^{\prime }\mathsf{z%
}^{N-1}\right) T\left( \mathsf{z}_{1}\mathsf{z}^{N-1}\right) T\left( \mathsf{%
z}_{1}^{\prime }\mathsf{z}^{N-1}\right) ^{\ast }d\mathsf{z}^{N-1}\overset{%
a.e.}{=}0\qquad \forall \breve{X}\in KerL_{N}^{1}
\end{equation}

\begin{itemize}
\item The two density contraction
\end{itemize}

In order that $KerL_{N}^{2}\rightarrow KerL_{N}^{2}$ by $\mathfrak{S}%
_{L_{N}^{2}}$ we must have that%
\begin{equation}
\int X\left( \mathsf{z}_{1}\mathsf{z}_{2}\mathsf{z}^{N-2}\mathsf{,z}%
_{1}^{\prime }\mathsf{z}_{2}^{\prime }\mathsf{z}^{N-2}\right) T\left( 
\mathsf{z}_{1}\mathsf{z}_{2}\mathsf{z}^{N-2}\right) T\left( \mathsf{z}%
_{1}^{\prime }\mathsf{z}_{2}^{\prime }\mathsf{z}^{N-2}\right) ^{\ast }d%
\mathsf{z}^{N-2}\overset{a.e.}{=}0\qquad \forall \breve{X}\in KerL_{N}^{2}
\end{equation}

The semi groups $\mathfrak{S}_{\Xi }$, $\mathfrak{S}_{L_{N}^{1}}$ and $%
\mathfrak{S}_{L_{N}^{2}}$ are not algebras (they are not vector spaces).

\section{Approximation Theory}

\section{Outline of Constructive Approach to Functionals I: DFT}

\bigskip

\bigskip

\section{Outline of Constructive Approach to Functionals II: DMFT (FORDO\
Functionals)}

\bigskip

\bigskip

\section{\protect\bigskip Outline of Constructive Approach to Functionals
III: N-Representability Problem (SORDO\ Functionals)}

\end{document}
